% =====================================================================
%  FICHE 4 — THEME 1 — Exercices niveau MOYEN
% =====================================================================
\documentclass[11pt,a4paper]{article}
\usepackage[utf8]{inputenc}
\usepackage[T1]{fontenc}
\usepackage{lmodern}
\usepackage[french]{babel}
\usepackage{amsmath,amssymb,mathtools}
\usepackage{icomma}
\usepackage[version=4]{mhchem}
\usepackage{siunitx}
\sisetup{output-decimal-marker={,},per-mode=symbol}
\DeclareSIUnit\Molar{M}
\usepackage{xcolor}
\usepackage{tikz}
\usepackage[most]{tcolorbox}
\usepackage{enumitem}
\usepackage{booktabs}
\usepackage{array}
\usepackage{fancyhdr}
\usepackage[a4paper,margin=2.2cm,headheight=26pt,headsep=10pt]{geometry}

\definecolor{primary}{RGB}{146,95,160}
\definecolor{primarydark}{RGB}{92,52,104}
\definecolor{defcol}{RGB}{0,114,178}

\tcbset{boxbase/.style={enhanced,breakable,boxrule=0.9pt,arc=2.5pt,
   left=3mm,right=3mm,top=2mm,bottom=2mm,fonttitle=\bfseries,coltitle=white,
   toptitle=0.5mm,bottomtitle=0.5mm}}
\newtcolorbox{donnees}{boxbase,colback=defcol!5,colframe=defcol,title={\textsc{Données}}}

\newcounter{exo}
\newenvironment{exo}[1][]{%
  \par\medskip\refstepcounter{exo}%
  \noindent\begin{tcolorbox}[boxbase,colback=primary!7,colframe=primary,
     title={\textsc{Exercice \theexo}\ifx\relax#1\relax\relax\else\textmd{ --- \itshape #1}\fi}]}%
  {\end{tcolorbox}}

\pagestyle{fancy}\fancyhf{}
\fancyhead[L]{\small\textcolor{primarydark}{\textbf{Fiche 4} --- Th.\,1 : Mole, masse molaire, entit\'es}}
\fancyhead[R]{\small\textcolor{primarydark}{\textsc{Chimie} --- \textbf{Moyen}}}
\fancyfoot[C]{\small\thepage}
\renewcommand{\headrulewidth}{0.8pt}
\renewcommand{\headrule}{\hbox to\headwidth{\color{primary}\leaders\hrule height \headrulewidth\hfill}}
\newcommand{\NA}{N_{\!\text{A}}}
\setlength{\parindent}{0pt}

\begin{document}

\begin{center}
{\LARGE\bfseries\color{primarydark} Thème 1 --- Niveau Moyen}\\[2pt]
{\color{primary}\itshape Enchaîner les conversions, extraire une part}
\end{center}

\begin{donnees}
Masses molaires atomiques (\si{\gram\per\mole}) :
$\ce{H}=1$ ; $\ce{C}=12$ ; $\ce{N}=14$ ; $\ce{O}=16$ ; $\ce{Na}=23$ ; $\ce{S}=32$ ;
$\ce{Al}=27$ ; $\ce{Fe}=56$ ; $\ce{He}=4$.\quad
$\NA=\SI{6.02e23}{\per\mole}$.
\end{donnees}

\begin{exo}[chaîne masse $\to$ moles $\to$ molécules]
L'aspirine (acide acétylsalicylique) a pour formule brute \ce{C9H8O4}. Un comprimé contient
$\SI{360}{\milli\gram}$ d'aspirine.
\begin{enumerate}[label=\textbf{\alph*)},leftmargin=2em,topsep=2pt,itemsep=2pt]
  \item Calculer la masse molaire de l'aspirine.
  \item En déduire le nombre de \textbf{molécules} d'aspirine contenues dans un comprimé.
\end{enumerate}
\end{exo}

\begin{exo}[compter des ions à partir d'une masse]
Pendant la grossesse, des compléments alimentaires apportent du sulfate de fer(II) \ce{FeSO4}.
Un comprimé contient $\SI{304}{\milli\gram}$ de \ce{FeSO4}. Sachant que chaque unité \ce{FeSO4}
libère un ion \ce{Fe^2+}, déterminer le \textbf{nombre d'ions \ce{Fe^2+}} apportés par un comprimé.
\end{exo}

\begin{exo}[pourcentage massique d'un élément]
\begin{enumerate}[label=\textbf{\alph*)},leftmargin=2em,topsep=2pt,itemsep=2pt]
  \item L'urée a pour formule \ce{CH4N2O}. Calculer le \textbf{pourcentage massique de carbone}
        dans l'urée.
  \item L'hydrogénosulfite de sodium a pour formule \ce{NaHSO3}. Calculer le \textbf{pourcentage
        massique d'oxygène} dans cette molécule.
\end{enumerate}
\end{exo}

\begin{exo}[doses de micronutriments et conversions d'unités]
Les apports journaliers recommandés d'un adulte comprennent $\SI{14}{\milli\gram}$ de fer (atomes
\ce{Fe}) et $\SI{48.8}{\micro\gram}$ de vitamine B8, de formule \ce{C10H16N2O3S}.
\begin{enumerate}[label=\textbf{\alph*)},leftmargin=2em,topsep=2pt,itemsep=2pt]
  \item Calculer la quantité de matière (nombre de moles) d'atomes de fer.
  \item Calculer la masse molaire de la vitamine B8, puis la quantité de matière correspondante.
\end{enumerate}
\end{exo}

\begin{exo}[comparer des nombres d'atomes]
On dispose de trois échantillons : $\SI{2.0}{\gram}$ d'hélium \ce{He} (gaz monoatomique),
$\SI{8.0}{\gram}$ de dioxygène \ce{O2} et $\SI{9.0}{\gram}$ d'aluminium \ce{Al}.\\
Pour chaque échantillon, calculer le \textbf{nombre total d'atomes} qu'il contient, puis indiquer
lequel (ou lesquels) en contient le plus.
\end{exo}

\end{document}
